\section{\MSAGA{} Kernel Programming Model}

The AWS Neuron Kernel Interface (NKI)~\cite{aws_neuron_nki} provides an expressive 
Python interface that allows users to write custom kernels for AWS Neuron 
devices. It enables instruction-level control with the convenience of 
Python APIs.
Inspired by AWS NKI, we design a similar but simplified Python library 
for writing custom \MSAGA{} kernels. The library consists of three main 
components: type-safe tensors, a Python API for the \MSAGA{} instruction set, 
and a lightweight JIT compiler.

\subsection{Type-Safe Tensors}

The \MSAGA{} Python library implements a PyTorch-like~\cite{pytorch,pytorch2} tensor abstraction 
for representing multi-dimensional arrays. Because the memory hierarchy 
is fully user-managed, tensors can be allocated in one of three memory 
spaces: main memory, scratchpad SRAM, or accumulation SRAM. Accordingly, 
we define three tensor types:

\begin{itemize}
  \item \texttt{MTile} — tensor allocated in main memory,
  \item \texttt{STile} — tensor allocated in scratchpad SRAM,
  \item \texttt{ATile} — tensor allocated in accumulation SRAM.
\end{itemize}

All tensor types support a subset of the PyTorch API, including 
\texttt{shape}, \texttt{dtype}, \texttt{split}, and \texttt{reverse}. 
By explicitly distinguishing tensor types, users can write custom 
kernel functions that clearly declare the expected memory scope of input and output 
tensors.

\subsection{Python API for \MSAGA{} Instructions}

Each \MSAGA{} instruction is exposed through a corresponding Python API. 
These APIs are designed to be type-safe, ensuring correct usage of 
tensor types. The complete set of supported APIs is shown in 
\autoref{lst:msaga_api}.
Internally, these functions construct hardware instructions using the 
metadata of the input tensors, such as shape, stride, and data type. 
This abstraction allows users to write high-level kernels without 
managing low-level hardware details.

\begin{listing}
\begin{mintedpython}
# DMA Instructions
def store_tile(src: ATile, dst: MTile, ...)
def load_tile(src: MTile, dst: STile, ...)
# Compute Instructions
def load_stationary(tile: STile, ...)
def attn_score(K: STile, l: ATile, ...)
def attn_value(V: STile, O: ATile, ...)
def reciprocal(l: ATile, ...)
def attn_lse_norm(O: ATile, ...)
\end{mintedpython}
\caption{\MSAGA{} Python instructions.}
\label{lst:msaga_api}
\end{listing}

\subsection{JIT Kernel Compiler}

The JIT compiler takes a kernel function written with the \MSAGA{} Python 
library and compiles all internally constructed \MSAGA{} instructions into a 
binary program that can be directly submitted to the \MSAGA{} accelerator.
It processes Python functions decorated with \texttt{@F.kernel}, which also
allow users to specify the target device for execution: either a simulated device using 
Verilator~\cite{verilator} or a real device running on Xilinx Alveo U55C FPGA.
Device-host communication is also handled automatically by the compiler.


\autoref{lst:msaga_fa_kernel} shows a complete implementation of the \SA{} kernel using the Python library.
In this example, the device is simulated by Verilator, while the host compiles the entire kernel into a binary program,
loads input tensors into the main memory simulated by DRAMSim2~\cite{dramsim2},
submits the compiled program to the device, and retrieves the output tensors back from DRAMSim2 after the computation is done.
As \MSAGA{} currently does not support matrix transposition,
the host transposes the input matrix $V$ and the output matrix $O$.
On commercial hardware, such transpositions can be performed by the DMA engine 
during memory transfers~\cite{aws_neuron_nki}, and therefore should not incur 
significant overhead.

\begin{listing}
\begin{mintedpython}
import asa as F

@F.kernel(device="verilator_sim")
def attention(Q: MTile, K: MTile, Vt: MTile):
  # allocate output tensor
  Ot: MTile = F.alloc_mem(..)
  # split large tensors into tiles
  Ot_MTiles = Ot.split(br, dim=-1) # [d, br]
  Q_MTiles = Q.split(br, dim=-2)   # [br, d]
  K_MTiles = K.split(bc, dim=-2)   # [br, d]
  Vt_MTiles = Vt.split(bc, dim=-1) # [d, bc]
  # double buffering for Q, K, Vt
  K_STiles = (F.alloc_spad(..),
              F.alloc_spad(..))
  Vt_STiles = (F.alloc_spad(..),
               F.alloc_spad(..))
  # allocate space for accumulation results
  log_expsum = F.alloc_accum(..)   # [1, br]
  Ot_ATile = F.alloc_accum(..)     # [d, br]

  for i, Q_i in enumerate(Q_MTiles):
    F.load_tile(Q_i, Q_STiles[i % 2])
    for j, (K_j, Vt_j) in \
        enumerate(zip(K_MTiles, Vt_MTiles)):
      F.load_stationary(Q_STiles[i % 2])
      F.load_tile(K_j, K_STiles[j % 2])
      F.attn_score(K_STiles[j % 2], log_expsum)
      F.load_tile(Vt_j, Vt_STiles[j % 2])
      F.attn_value(Vt_STiles[j % 2], Ot_ATile)
    F.reciprocal(log_expsum)
    F.attn_lse_norm(Ot_ATile)
    F.store_tile(Ot_ATile, Ot_MTiles[i])

# Run simulation and return result
Ot = attention(Q, K, Vt)
O = Ot.to_numpy().T  # Host-side transpose
\end{mintedpython}
\caption{\MSAGA{} \FA{} kernel.}
\label{lst:msaga_fa_kernel}
\end{listing}

% The JIT compiler processes Python functions decorated with 
% \texttt{@F.kernel}. The decorator also specifies the target device 
% (e.g., a Verilator-based simulator~\cite{verilator}), and handles device-host 
% communication automatically.

% In the example above, the host loads input tensors into device memory using 
% DRAMSim2~\cite{dramsim2}, dispatches the compiled instructions to the \MSAGA{} device, 
% and retrieves the output tensors upon completion. 
% Because \MSAGA{} does not currently support hardware matrix transposition, 
% the input matrix $V$ is transposed in advance, and the output matrix $O$ 
% is converted to NumPy~\cite{numpy} arrays for post-processing. 
% On commercial hardware, such transpositions can be performed by the DMA engine 
% during memory transfers~\cite{aws_neuron_nki}, and therefore should not incur 
% significant overhead.
